\section{Introduction + Requirements}
\begin{itemize}
	\item shot noise requirement
	\item radiation pressure noise setting noise requirement
	\item inner+outer loop design
	\item could include calculation for mismatches in the IFO? Point absorbers etc. 
	\item Patrick Kwee's Thesis. 
	\item Useful ligo reference \href{https://dcc.ligo.org/cgi-bin/DocDB/ShowDocument?docid=T010075}{T010075}
\end{itemize}
\subsection{Intensity Noise}
\href{https://alog.ligo-wa.caltech.edu/aLOG/uploads/46276_20190107234129_SRCLArmPowerMeas.pdf}{SRC dither measurement explanation}.\href{https://alog.ligo-la.caltech.edu/aLOG/index.php}{Measurement 12/11/2019}. Amount of imbalance  $\sim 3\%$ \href{https://alog.ligo-la.caltech.edu/aLOG/index.php?callRep=43298}{Terra Intensity noise injection method}. LIGO-T010075-v3 states 1\% arm imbalance for nominal value. LIGO Document T050036-v4 also states 1\%, this is probably a better reference.
\subsection{Coupling of Intensity Noise}
In the GW band, laser intensity noise couples via differential arm finesse (i.e. losses) -> difference in arm powers -> radiation pressure fluctuations which look like gw signal. Would be nice to have an equation.
\paragraph{}
Power fluctuations in LO light -> can't be distinguished from GW signal. Will be cancelled with BHD.
\paragraph{}
PRC passively filters intensity noise of laser. Bandwidth $\approx 1\si{\hertz}$.
\subsection{Control band fluctations 0.1-10Hz}
Creates radiation pressure -> GW signal coupling. 
\subsection{RF noise}
RF signals used to control the IFO -> Noise at these frequencies degrades SNR -> indirect coupling of noise to GW channel. PMC used for this. IMC cannot be used as EOM before this. NPRO is dominant RF noise source.
\subsection{Inner Loop}

\subsection{Outer Loop}

\section{Free Running Noise and Noise Hunting}
\begin{itemize}
	\item free running noise of the laser
	\item what dominates the laser noise? NPRO
	\item noise sources such as mounts 
	\item noise from PMC offset not being tuned
	\item 4Hz noise from Beckhoff
\end{itemize}
\section{AOM Stabilisation}
\begin{itemize}
	\item description of AOM
	\item AOM TF
	\item AOM loop TF
	\item plot of AOM stabilisation working
\end{itemize}
\begin{figure}
	\centering
	\includegraphics[width=0.7\linewidth]{figures/AOM_stabilisation_noise}
	\caption{Free running and stabilised noise of laser}
	\label{fig:aomstabilisationnoise}
\end{figure}
\begin{figure}
	\centering
	\includegraphics[width=0.7\linewidth]{figures/AOM_stabilisation_TF}
	\caption{Transfer functions for parts of the AOM servo}
	\label{fig:aomstabilisationtf}
\end{figure}
\section{Current Shunt Stabilisation}
\begin{itemize}
	\item description of CS
	\item CS TF
	\item CS loop TF
	\item plot of CS stabilisation working
\end{itemize}
\begin{figure}
	\centering
	\includegraphics[width=0.7\linewidth]{figures/current_shunt_tfs}
	\caption{current shunt input to light power transfer function for different diode currents}
	\label{fig:currentshunttfs}
\end{figure}
\begin{figure}
	\centering
	\includegraphics[width=0.7\linewidth]{figures/dc_power_all_pd}
	\caption{DC behaviour of current shunt}
	\label{fig:dcpowerallpd}
\end{figure}
\section{Comparison of Two stabilisation Schemes}
\begin{itemize}
	\item why I think the AOM is a better choice even though they both work
\end{itemize}
\section{Conclusion}